% Desenvolvido por André Leite - DE - UFPE
% Ajustes Raydonal Ospina - DE - UFPE/UFBA

\documentclass[a4paper,12pt,oneside,twocolumn]{Config/milktest}

%---------------------- customização ----------------------



\usepackage{fontspec}
\usepackage{xunicode}
\usepackage{xltxtra}
\usepackage[LGR,T1]{fontenc}
%\usepackage{ebgaramond}

%\usepackage[cp1252]{inputenc}                %% Use um dos dois esquemas
\usepackage[english,brazil]{babel}                 %% Hifenizacao e estilo
\usepackage{pstricks}                          %% Pacote de comandos p/ Facilitar o Postscripts
\usepackage{pstricks-add}                    %% " "
\usepackage{pst-all,pst-math,pst-plot, pst-func}       %% " "
\usepackage{textcomp}							%%
%\usepackage[tight,nice]{units}				%%
\usepackage[centertags]{amsmath}       %% [leqno]
\usepackage{amsfonts}                         %% Pacotes da AMS
\usepackage{amssymb}                         %%
\usepackage{stmaryrd} 							%%
\usepackage{Config/milkstyle}
%\usepackage{sabon}

\usepackage{nicefrac}%  Nice Frescura
\usepackage{balance}%
\usepackage{lastpage} %                      %% Referênciar a última página
\usepackage{graphicx}                        %% Mais ferramentas para
%\usepackage[T1]{pbsi} % Procurar e instalar no Mac!!!!! <------<-------<------
\usepackage{mathrsfs}                        %% $\mathscr{ABCD}$
\usepackage{lettrine}                        %% ...
%\usepackage[normal,bf,up]{caption2}          %% Controlar melhor os captions....
\usepackage{eucal}                           %% beau
\usepackage{colortbl} %
\usepackage{hyperref}                 %%
%\usepackage{amsthm}
\usepackage{indentfirst}%
\usepackage{enumitem}%
\usepackage{colordvi}                        %% Usar Cores...
\usepackage{fancyhdr}%
\usepackage{yfonts}%
\usepackage{theorem}%
\usepackage{multirow}                        %% Juntar linhas em tabelas
\usepackage{midfloat}%
\usepackage{cancel}%
%\usepackage[square]{natbib}%sectionbib,
\usepackage{fancyhdr} % cabeçalhos ...
\usepackage{rotating}


\usepackage{lipsum} 
%% Lista de Pacotes utilizados


%% Deixando o caption mais interessante...
%\renewcommand{\captionfont}{\itshape}


\hypersetup{%backref,%
%            pdfpagemode=FullScreen,%
            colorlinks=true,%
            citecolor=blue,%
            urlcolor=blue, %
            % urlcolor=black, %
            linkcolor=blue,%
            pdfauthor={André Leite Wanderley}, %
            pdftitle={1EE/ET406}, %
            pdfcreator={André Leite Wanderley}, %
            breaklinks=true, %
            linktocpage=true
            }%


\DeclareMathOperator*{\Inf}{Inf}%
\DeclareMathOperator*{\Sup}{Sup}%
\DeclareMathOperator*{\Min}{Min}%
\DeclareMathOperator*{\Max}{Max}
\DeclareMathOperator*{\Sgn}{Sgn}%
\DeclareMathOperator*{\E}{E}%
\newcommand{\mbf}{\boldsymbol}%
\DeclareMathAlphabet\mathbb{U}{fplmbb}{m}{n}%
%
%\numberwithin{equation}{section}                       %%
%\numberwithin{table}{section}                          %% Numeração dos floats...
%\numberwithin{figure}{section}                         %%


\newcommand{\oo}{\mathbin{\bullet}}
\newcommand{\oc}{\mathbin{\circ}}
\newcommand{\noo}{\mathbin{\not\!\bullet}}
\newcommand{\cf}{\emph{cf.}\ }
\newcommand{\ie}{\emph{i.e.}}
\newcommand{\eg}{\emph{e.g.}}
\newcommand{\raab}{\emph{reductio ad absurdum}}
\newcommand{\eetc}{\textit{et c\ae tera}}
\newcommand{\dom}{\mathfrak{D}}
\newcommand{\ima}{\mathfrak{I}}
\newcommand{\idenA}{\mathcal{I}_\mathcal{A}}
\newcommand{\ce}{\ensuremath{\mathcal{X}}}
\newcommand{\piq}{\ensuremath{\succsim}}
\newcommand{\pq}{\ensuremath{\succ}}
\newcommand{\iq}{\ensuremath{\sim}}
\newcommand{\niq}{\ensuremath{\nsim}}
\newcommand{\npiq}{\ensuremath{\not\succsim}}
\newcommand{\npq}{\ensuremath{\nsucc}}
\newcommand{\po}[2]{\ensuremath{(#1,#2)}}

%\setlist{noitemsep}
%\setlist{labelindent=\parindent} % << Usually a good idea
\setlist{itemsep=0pt}%
\setitemize{leftmargin=\parindent} %
\setitemize[1]{label=$\triangleright$} %
%\setenumerate[1]{label=\arabic*.,ref=\arabic*}
\setenumerate[1]{label=\arabic*.,ref=\arabic*}
\setenumerate[2]{label=\theenumi.\arabic*.,ref=\theenumi.\arabic*}
\setenumerate[3]{label=\roman*.,ref=\theenumii.\roman*}
%\setdescription{%font=\sffamily\bfseries}

%                }%


%%%%%%%%%%%%%%%%%%%%%%%%%%%%%%%%
% Para facilitar minha vida

\newcommand{\Ds}{\displaystyle}%
\newcommand{\Ts}{\textstyle}%
\newcommand{\Ss}{\scriptstyle}%
\newcommand{\SSs}{\scriptscriptstyle}%

\def\aureos{%
            \newrgbcolor{gray1}{.9 .9 .9}
            \newrgbcolor{guess1}{1 0.388235294 0.090196078}
            \psset{unit=1in,cornersize=absolute}%
            \psline[linecolor=guess1,linewidth=1.2pt](-0.125,-1.291667)(1.1125,-1.291667)
            \psline[linecolor=guess1,linewidth=0.6pt](0.347683,-1.291667)(0.347683,-0.52685)
            \psline[linecolor=guess1,linewidth=0.6pt](0.347683,-0.999533)(-0.125,-0.999533)
            \psline[linecolor=guess1,linewidth=0.6pt](0.167134,-0.999533)(0.167134,-1.291667)
            \psline[linecolor=guess1,linewidth=0.6pt](0.167134,-1.111118)(0.347683,-1.111118)
            \psline[linecolor=guess1,linewidth=0.6pt](0.236098,-1.111118)(0.236098,-0.999533)
            \psline[linecolor=guess1,linewidth=0.6pt](0.236098,-1.068496)(0.167134,-1.068496)
            \psline[linecolor=guess1,linewidth=0.6pt](0.209756,-1.068496)(0.209756,-1.111118)
            \psline[linecolor=guess1,linewidth=0.6pt](0.209756,-1.084776)(0.236098,-1.084776)
           }%

\def\GuEsS{%
            \newrgbcolor{guess1}{1 0.388235294 0.090196078}
            \newrgbcolor{guess2}{0.960784314 0.109803922 0.043137255}
            \psset{unit=1in,cornersize=absolute}%
            \rput(0,-1.5){{\fontsize{35}{0}\usefont{OT1}{cmr}{m}{n}\selectfont\textsc{\textbf{S}}}}%
            \rput(0.625,-1.5){\shortstack{{\huge\textbf{\textsc{ystems}}}\\
                {\scriptsize\sc\hfill Research Group}}}
            \psline[linecolor=black,fillcolor=guess2,fillstyle=none,linewidth=3.6pt](-0.125,-1.725)(1.1125,-1.725)%
            \uput{0.5ex}[l](1.148214,-1.229167){\llap{{\tiny\sc UFBA}}}%
            \psline[linecolor=guess1,fillcolor=guess1,fillstyle=none,linewidth=1.2pt](-0.125,-1.291667)(1.1125,-1.291667)%
            \psline[linecolor=guess1,fillcolor=guess1,fillstyle=none,linewidth=.6pt](0.347683,-1.291667)(0.347683,-0.52685)%
            \psline[linecolor=guess1,fillcolor=guess1,fillstyle=none,linewidth=.6pt](0.347683,-0.999533)(-0.125,-0.999533)%
            \psline[linecolor=guess1,fillcolor=guess1,fillstyle=none,linewidth=.6pt](0.167134,-0.999533)(0.167134,-1.291667)%
            \psline[linecolor=guess1,fillcolor=guess1,fillstyle=none,linewidth=.6pt](0.167134,-1.111118)(0.347683,-1.111118)%
            \psline[linecolor=guess1,fillcolor=guess1,fillstyle=none,linewidth=.6pt](0.236098,-1.111118)(0.236098,-0.999533)%
            \psline[linecolor=guess1,fillcolor=guess1,fillstyle=none,linewidth=.6pt](0.236098,-1.068496)(0.167134,-1.068496)%
            \psline[linecolor=guess1,fillcolor=guess1,fillstyle=none,linewidth=.6pt](0.209756,-1.068496)(0.209756,-1.111118)%
            \psline[linecolor=guess1,fillcolor=guess1,fillstyle=none,linewidth=.6pt](0.209756,-1.084776)(0.236098,-1.084776)%
           }%

%%%%%%%%%%%%%%%%%%%%%%%%%%%%%%%%%%%%%%%%%%%%%%%%%%%%%%%%


\theoremstyle{plain}                     %%
%\theoremstyle{margin}                   %%
\theoremheaderfont{\scshape\bfseries}    %%
\theorembodyfont{\rmfamily}                       %%

\newtheorem{Prop}{\small Proposition}[section]  %%
\newtheorem{Def}{\small Definition}[section]    %%
\newtheorem{Not}{\small Notice}[section]    %%
\newtheorem{Theo}{\small Theorem}[section]      %%
\newtheorem{Lem}{\small Lema}[section]         %%
\newtheorem{Coro}{\small Corollary}[section]    %%
\newtheorem{Conj}{\small Conjecture}[section]   %%
\newtheorem{Prob}{\small Problem}[section]      %%
\newtheorem{Exa}{\small Example}[section]       %%
%\theoremheaderfont{\bfseries}    %%
%\theorembodyfont{\rmfamily}  %%
\newtheorem{proof}{\small Demonstration}[section]
\newtheorem{sol}{\small Solution}[section]


\newenvironment{Proofm}               % Deve ter uma maneira
    {%\strut \par\medskip             % ridiculamente mais
     %\strut\hspace{\parindent}      % fácil de fazer isso...
     %\begin{minipage}{.8\linewidth} % I'm sure about that.
     \begin{proof}\small}                   %
    {\strut\hfill\textsc{q.e.d.}\end{proof}                      %
    %\end{minipage}                  %
    %\strut\par\medskip
                        }             %

\newenvironment{solm}               % Deve ter uma maneira
    {%\strut\par%\medskip             % ridiculamente mais
     %\strut\hspace{\parindent}      % fácil de fazer isso...
     %\begin{minipage}{.8\linewidth} % I'm sure about that.
     %\begin{sol}
     \small
     \noindent\textbf{Solução:}
     %\par %\smallskip
     }                   %
    {\strut\hfill\textsc{$\blacksquare$}%\end{sol}                      %
    %\end{minipage}                  %
    %\strut\par\medskip
                        }             %



\newenvironment{Propm}               % Deve ter uma maneira
    {\strut\par\smallskip             % ridiculamente mais
     \strut\hspace{\parindent}      % fácil de fazer isso...
     \begin{minipage}{.8\linewidth} % I'm sure about that.
     \begin{Prop}\small}                   %
    {\end{Prop}                      %
    \end{minipage}                  %
    \strut\par\medskip
    }             %


\newenvironment{Defm}               % Deve ter uma maneira
    {\strut\par\smallskip             % ridiculamente mais
     \strut\hspace{\parindent}      % fácil de fazer isso...
     \begin{minipage}{.8\linewidth} % I'm sure about that.
     \begin{Def}\small}                   %
    {\end{Def}                      %
    \end{minipage}                  %
    \strut\par\medskip}             %

\newenvironment{Theom}               % Deve ter uma maneira
    {\strut\par\smallskip             % ridiculamente mais
     \strut\hspace{\parindent}      % fácil de fazer isso...
     \begin{minipage}{.8\linewidth} % I'm sure about that.
     \begin{Theo}\small}                   %
    {\end{Theo}                      %
    \end{minipage}                  %
    \strut\par\medskip}

\newenvironment{Notm}               % Deve ter uma maneira
    {\strut\par\smallskip             % ridiculamente mais
     \strut\hspace{\parindent}      % fácil de fazer isso...
     \begin{minipage}{.8\linewidth} % I'm sure about that.
     \begin{Not}\small}                   %
    {\end{Not}                      %
    \end{minipage}                  %
    \strut\par\medskip}


\newenvironment{Exam}               % Deve ter uma maneira
    {%\strut\par\smallskip             % ridiculamente mais
     %\strut\hspace{\parindent}      % fácil de fazer isso...
     %\begin{minipage}{1\linewidth} % I'm sure about that.
     \begin{Exa}\small}                   %
    { \scriptsize$\blacksquare$\end{Exa}                      %
    %\end{minipage}                  %
    \strut\par\medskip}

\def\Pt#1(#2,#3){
       \expandafter\def\csname #1\endcsname{#2,#3}%
       \expandafter\def\csname #1X\endcsname{#2}%
       \expandafter\def\csname #1Y\endcsname{#3}}%

\newcommand{\hlinem}[3]{\noalign{\vspace{#3}\hrule height #1 depth #2}}
\newcommand{\hlinemilk}[3]{\noalign{}}

\newcommand{\quem}[1]{\par\medskip\strut\hfill\bf  --- #1}
\newcounter{questionNumber}
\newcommand{\question}{\stepcounter{questionNumber}\lettrine[lhang=.75, findent=10pt, nindent=0pt,lines=3, lraise=.2]{\thequestionNumber}{$\blacktriangleright$}}


\newenvironment{milkquote}
    {\medskip %
    \begin{center}
    \begin{minipage}{.9\linewidth}
    \hspace{.25cm}\small\itshape}%
    {\medskip\end{minipage}
    \end{center}
    }%


\newcommand{\fim}{%
                 \par\medskip
                 \begin{center}
                      \includegraphics[width=1cm]{ornaments/ornamento4.eps}
                 \end{center}
                 \vspace*{.5cm}
                 }%



\newenvironment{milkanswer}
    {\par %
    \begin{center}
    \color{orange}\textbullet\quad\textbf{Solution}\quad\textbullet
    \end{center}\small
    }%
    {
    \fim
    }


%%%
%%% Set up page layout parameters.
%%%
\setlength\paperheight{297mm}      %
\setlength\paperwidth{210mm}       % papel a4
%\setlength{\parindent}{6mm}

\addtolength{\hoffset}{-1in}%
\addtolength{\voffset}{-1.5in}

\setlength{\oddsidemargin}{25mm}    % 
\setlength{\topmargin}{5cm}        %
\setlength{\textwidth}{170mm}      % Margem direita: 25mm
\setlength{\textheight}{220mm}  % Margem Inferior: 25mm
\setlength{\marginparsep}{0mm}     %

\setlength{\headheight}{16mm}
\renewcommand{\columnsep}{1cm}%
%\marginsize{1.25in}{1.75in}{1.25in}{1.25in}



% If you can to leave the header here
%---

%---

\renewcommand{\headrulewidth}{0pt}
\definecolor{cor}{rgb}{.5,0, 0}

%%%%%%%%%%%%%%%%%%%%%%%%%%% Font 

\defaultfontfeatures{Ligatures=Historic,Contextuals=Alternate, Numbers=OldStyle,RawFeature={ +hlig, +dlig, +ss05}}
\setmainfont[RawFeature={+ss05, +hlig, +dlig}, BoldFont = Gentium Basic Bold, ItalicFeatures={RawFeature={+clig},CharacterVariant=5:0}]{EB Garamond}
% % Do the replacements manually in titles, to not ruin the textsc
% \newfontfamily\booktitlefont[RawFeature={-ss02},LetterSpace=40,WordSpace=6]{EB Garamond}
% \newfontfamily\spacedfont[RawFeature={-ss02},LetterSpace=20,WordSpace=3]{EB Garamond}
% \newfontfamily\lettrinefont{EB Garamond Initials}
% \newfontfamily\headerfont{EB Garamond}

%\setmainfont[BoldFont = Gentium Basic Bold,
%      ItalicFont = Gentium Plus Italic,
%      BoldItalicFont = Gentium Basic Bold Italic,
%      Mapping = tex-text]{Gentium Plus}





 

\usepackage{bm}
\usepackage{verbatimbox}
\usepackage{float}

%%%%%%%%%%%%%----- Header------%%%%%%%%%%%%%%%%%%%
\newsavebox{\mygraphic}
\sbox{\mygraphic}{\includegraphics[scale=1.2]{Figuras/ufba.eps}}

\fancypagestyle{milktablestyle}{
    \fancyhead{}
    \fancyfoot[C]{\thepage/\pageref{LastPage}} %
}

\pagestyle{fancy}

\fancyhead{}
\fancyhead[C]{
\rput(0,2){
  \begin{minipage}{1.75cm}\centering
    \includegraphics[width=1.75cm]{Figuras/ufbaBlack.eps}
  \end{minipage}~\hfill~\begin{minipage}{10cm}\Large\centering
   \color{cor}{\scshape universidade federal da bahia}\\
    {\scshape instituto de matemática e estatística}\\
    {\scshape departamento de estatística}\\[5pt]
    % Aqui as informações da prova e a data
    {\Large\scshape  prova II --- matd44\\ } 
  \end{minipage}~\hfill~\begin{minipage}{1.75cm}\centering
    \includegraphics[width=1.75cm]{Figuras/ufba.eps}
  \end{minipage}
  }}

\fancyfoot[C]{\thepage/\pageref{LastPage}} % 

%%%%%%%%%%%%%----- Header------%%%%%%%%%%%%%%%%%%%

% \usepackage[printwatermark]{xwatermark}
% \newwatermark[allpages,color=gray!20,angle=45,scale=3,xpos=30,ypos=-60]{DRAFT}

\begin{document}




% %------------------------------------------------
\beb
{\small
\begin{description}
\item[Professor:] Raydonal Ospina
\item[E-mail:] raydonal@castlab.org
% \item[Site do curso:] \url{https://castlab.org/courses/matd44/}

\item[Regras:] Leia com atenção as perguntas. {\bf Todas as questões devem ser detalhadas.} A prova deve ser claramente resolvida. Seja claro e organizado. 

\item[Atenção:] Você deverá encaminhar a solução da sua prova digitalizada em formato {\bf PDF} no e-mail acima no dia 20/08/2023 até as 15:00h (GMT-3 - Horário de Brasília) 
\end{description}
}
\eeb

% %%%%%%%%%%%%%%%%%%%%%%%%%%%%%%%%%%%%%%%%%%%%%%%%%%%%%%%%%%%%%%%%%
\balance
%%%%%%%%%%%%%%%%%%%%%%%%%%%%%%%%%%%%%%%%%%%%%%%%%%%%%%%%%%%%%%%%%


% \medskip 
% \question~({\bf Amostragem Bernoulli}) Seja $s$ uma amostra obtida de um desenho amostral de tipo Bernoulli com probabilidades de inclusão $\pi_k = \pi$ para todo $k \in U$ (população). Seja $n_s$ o tamanho de amostra da amostra $s.$ Mostre que a probabilidade condicional de se obter $s$ dado $n_s$ é a mesma que a probabilidade obtida por uma amostragem aleatória simples sem substituição de tamanho fixado $n_s$ de $N$ (Tamanho da população). 


% \question~({\bf Amostragem Bernoulli}) Na elaboração de um estudo de rede de escritório, o seguinte esquema de amostragem sequencial foi proposto para selecionar uma amostra aleatória de dois intervalos de horários de escritório não adjacentes dentre os cinco intervalos 9-10, 10-11, 11-12, 12-13, 13-14 (rotulados de 1 a 5). Seleciona-se o primeiro intervalo de uma hora para a amostra com igual probabilidade a partir dos cinco intervalos. Em seguida, seleciona-se o segundo intervalo de uma hora, sem reposição, com igual probabilidade dos intervalos que não são adjacentes no tempo ao primeiro intervalo selecionado no primeiro sorteio.
% \begin{itemize}
% \item[a)] Determine as probabilidades de inclusão de primeira ordem. 
% \item[b)] Determine as probabilidades de inclusão de segunda ordem e determine se esse plano amostral induzido pelo esquema de amostragem proposto é mensurável? 
% \item[c)] Determine as covariâncias dos indicadores de pertencimento à amostra.
% \end{itemize}


\medskip 
\question~~({\bf Amostragem estratificada}) A Tabela \ref{tab:2} contem as informações do gasto mensal em serviços públicos de uma amostra aleatória estratificada de 120 famílias na cidade de Salvador a qual foi geograficamente dividida em três estratos: Norte, Centro e Sul.

\begin{table}[!htb]
\begin{tabular}{l|ccc}
             & \multicolumn{3}{c}{Estratos}   \\ \hline
Estatísticas & Norte (1)   & Centro (2) & Sul (3)      \\ \hline
$N_h$        & 4000     & 6000     & 10000    \\
$W_h$        & 0,3      & 0,2      & 0,5      \\
$n_h$        & 40       & 36       & 44       \\
$\bar{y}_h$  & 1,2      & 2,4      & 0,6      \\
$\bar{Y}_h$  & 9600     & 7200     & 6000     \\
$s^2_h$             & 0,36     & 1,21     & 0,04     \\
${\rm Var}(\bar{y}_h)$          & 0,000993 & 0,004404 & 0,000226 \\ \hline
\end{tabular}
\caption{Informações do gasto familiar mensal em serviços públicos (em salários mínimos) a partir de uma amostra aleatória simples estratificada na cidade de Salvador.}
\label{tab:2}
\end{table}

\begin{itemize}
\item[a)]  Estime o gasto médio e o gasto total de toda a população. Estabeleça um intervalo de confiança de 95\% para essas medidas.
\item[b)] Suponhamos que os custos de coletar a informação por família para cada um dos estratos são: $C_1={\rm R}\$5.000,$ $C_2={\rm R}\$3.000$
 e $C_3={\rm R}\$1.000,$ respectivamente. Assuma que que existem informações previas em relação a variabilidade de cada estrato e que estas informações correspondem as estimativas $s^2_h$ da tabela \ref{tab:2}. Se o orçamento da coleta de informação não pode ser superior a ${\rm R}\$2,5$ milhões qual deve ser o tamanho da amostra global e como ele deve ser dividido entre os diferentes estratos? Suponha que a função de custo é linear na forma
 $$
 C = C_0+\sum_{h=1}^HC_h n_h,$$ em que $C$ 
 é o orçamento total para a coleta de informação, $C_0$ é o custo fixo que não depende do número de unidades amostrais a serem selecionadas e $C_h$ é o custo de amostrar uma unidade amostral no estrato $h.$ 
 
 \end{itemize}

\medskip 
\question~({\bf Estimação em subpopulações finitas (domínios)}) Uma amostra aleatória simples e sem substituição de 56 pessoas foi selecionada de uma população de 1000 trabalhadores da empresa LINCATECH. Foram coletadas informações sobre a renda mensal em miles de reais (Renda) e o sexo (Sexo) do trabalhador. Com as informações da Tabela \ref{ID:Renda} estime:

\begin{itemize}
%     \item[a)] A renda média dos trabalhadores. Estabeleça um intervalo de confiança de 95\% para a renda média.
% \item[b)] A renda total dos trabalhadores. Estabeleça um intervalo de confiança de 95\% para a renda total. 
\item[a)] Estime a proporção e o número total de mulheres na empresa. Estabeleça intervalos de confiança de 95\% para a proporção e total de mulheres na empresa.
\item[b)] Podem ser consideradas válidas aproximações pela distribuição normal no item anterior? Explique.
\item[c)] Considera que as amostras, tanto de homens como de mulheres poderiam ser assumidas como amostras aleatórias simples e sem substituição das respectivas subpopulações de homens e mulheres da empresa? Explique.
\item[d)] Como poderia ser estimada a renda média e o total das mulheres para toda a empresa se não se conhecesse  o número total delas?
\item[e)] Qual das duas subpopulações (homens, mulheres) é mais homogênea em relação a renda?
\end{itemize}

{\small \begin{table}[!htp]
\centering
\begin{tabular}{rlr|rlr}
  \hline
ID & Sexo & Renda & ID & Sexo & Renda \\ 
  \hline
    1 & Fem & 2094,90 &   15 & Mas & 2939,21 \\
    2 & Mas & 2386,14 &   16 & Mas & 1722,62 \\
    3 & Mas & 1562,82 &   17 & Fem & 2739,79 \\
    4 & Mas & 1781,41 &   18 & Mas & 1821,61 \\
    5 & Mas & 1603,14 &   19 & Mas & 1742,40 \\
    6 & Mas &  479,73 &   20 & Mas & 1845,22 \\
    7 & Mas & 2196,85 &   21 & Mas & 1916,60 \\
    8 & Mas & 2365,21 &   22 & Mas & 1329,28 \\
    9 & Mas & 2016,32 &   23 & Mas & 2143,04 \\
   10 & Mas & 1322,23 &   24 & Mas & 2618,97 \\
   11 & Fem & 2589,08 &   25 & Fem & 1399,85 \\
   12 & Fem & 2896,28 &   26 & Mas & 1610,41 \\
   13 & Mas & 1370,55 &   27 & Mas & 2300,84 \\
   14 & Mas &  975,94 &   28 & Mas & 1192,90 \\ \hline   
   29 & Mas & 2715,62 &   43 & Mas & 1802,16 \\
   30 & Mas & 2042,58 &   44 & Mas & 2444,68 \\
   31 & Mas & 2235,73 &   45 & Mas & 2644,75 \\
   32 & Mas & 2223,33 &   46 & Fem & 1431,53 \\
   33 & Mas & 2618,16 &   47 & Mas & 1094,02 \\
   34 & Mas & 2206,57 &   48 & Mas & 1548,96 \\
   35 & Mas & 2432,08 &   49 & Mas & 2410,51 \\
   36 & Mas & 1340,94 &   50 & Mas & 2286,37 \\
   37 & Mas & 2321,37 &   51 & Mas & 1589,07 \\
   38 & Mas & 1922,91 &   52 & Fem & 1646,21 \\
   39 & Mas & 2520,83 &   53 & Mas & 3358,63 \\
   40 & Mas & 2063,78 &   54 & Mas & 1369,42 \\
   41 & Mas & 2335,66 &   55 & Mas & 2047,63 \\
   42 & Mas & 2357,94 &   56 & Mas & 1719,34 \\
   \hline
\end{tabular}
\caption{Tabela de Informações dos empregados na amostra} 
\label{ID:Renda}
\end{table}
}




\medskip 
\question~({\bf Estimação por conglomerados}) suponha que o objetivo de uma pesquisa seja estimar a renda média em um bairro da cidade. Assuma que neste bairro existam \(N_I = 60\) quadras. {\scriptsize (A quadra ou quarteirão é a menor área de espaço urbano delimitada por ruas, rios ou avenidas, caracterizando a unidade básica de formação destes espaços)} 

É realizado um plano de amostragem aleatória simples de conglomerados e são selecionadas \(n_I = 5\) quadras, nas quais todos os domicílios são entrevistados. Os resultados da pesquisa são apresentados na Tabela \ref{tabela2} a seguir:

{\small \begin{table}[!htp]
\centering

\begin{tabular}{|p{2cm}|p{2cm}|p{2cm}|}
\hline
\textbf{ID da Quadra} & \textbf{Domicílios na Quadra} & \textbf{Renda Total na Quadra} \\ \hline
AW45 & 120 & 25000 \\ \hline
AW02 & 100 & 24000 \\ \hline
AW31 & 80  & 19000 \\ \hline
AW28 & 95  & 20100 \\ \hline
AW44 & 80  & 18000 \\ \hline
\end{tabular}
\caption{\label{tabela2}Tabela das cinco quadras selecionadas}
\end{table}
}
\begin{itemize}
    \item[(a)] Estime a renda total dos domicílios no bairro. Reporte o coeficiente de variação estimado.
    \item[(b)] Estime o número de domicílios no bairro. Reporte o coeficiente de variação estimado.
    \item[(c)] Assumindo que no bairro existam \(N = 2000\) domicílios, estime a renda média dos domicílios no bairro. Reporte o coeficiente de variação estimado.
    \item[(d)] Estime a renda média utilizando o estimador de Hájek\footnote{Lembrete: O estimador de Hájek (1971) é uma forma aprimorada do estimador de razão. O estimador de Hájek para a média populacional pode ser expresso como: \[
\widehat{\overline{Y}}_{\text{H}} = \frac{\sum_{i \in s} y_i}{\sum_{i \in s} \frac{1}{\pi_i}}.
\] O denominador é um estimador do tamanho populacional.}. Explique a diferença em relação à estimativa do ponto anterior. 
    
\end{itemize}

\bigskip


\centering\color{cor}{\scshape boa prova}



\end{document}

